\documentclass[a4paper]{article}

\usepackage[fontset=fandol]{ctex}
\usepackage{amsmath, amssymb}
\usepackage{graphicx}
\usepackage[margin=2.5cm]{geometry}
\usepackage[most]{tcolorbox}
\usepackage{etoolbox}
\usepackage{listings}
\usepackage{hyperref}
\usepackage{booktabs}
\usepackage{subcaption}
\usepackage{float}
\usepackage{tikz}
\IfFileExists{pgfplots.sty}{
    \usepackage{pgfplots}
    \pgfplotsset{compat=1.18}
}{}

\newtcolorbox{knowledgebox}[1]{
    enhanced, colback=blue!5!white, colframe=blue!75!black, colbacktitle=blue!75!black,
    coltitle=white, fonttitle=\bfseries, title=#1, attach boxed title to top left={yshift=-2mm, xshift=2mm},
    boxrule=1pt, sharp corners
}

\newtcolorbox{importantbox}[1]{
    enhanced, colback=yellow!10!white, colframe=yellow!80!black, colbacktitle=yellow!80!black,
    coltitle=black, fonttitle=\bfseries, title=#1, sharp corners
}

\newtcolorbox{warningbox}[1]{
    enhanced, colback=red!5!white, colframe=red!75!black, colbacktitle=red!75!black,
    coltitle=white, fonttitle=\bfseries, title=#1, sharp corners
}

\lstset{
    language=Python,
    basicstyle=\ttfamily\small,
    keywordstyle=\color{blue},
    stringstyle=\color{red!60!black},
    commentstyle=\color{green!60!black},
    breaklines=true,
    frame=single,
    numbers=left,
    numberstyle=\tiny\color{gray},
    captionpos=b,
    extendedchars=false
}

\newcommand{\notetitle}{CS294-277 课程笔记：Lecture 12 长程规划与语言}
\newcommand{\noteauthors}{Codex}
\newcommand{\notedate}{2026-04-16}
\newcommand{\videochannel}{CS294-277 / UC Berkeley}
\newcommand{\videopublishdate}{2024-12-02}
\newcommand{\videoduration}{未在本地元数据中提供}
\newcommand{\videourl}{https://www.youtube.com/watch?v=Ty_XNoXlQ5c}
\newcommand{\videocoverpath}{cover.jpg}

\begin{document}

\begin{titlepage}
\centering
{\Large 课程笔记\par}
\vspace{1.2cm}
{\huge\bfseries \notetitle\par}
\vspace{0.8cm}
{\large \noteauthors\par}
\vspace{0.3cm}
{\large \notedate\par}
\vspace{1.2cm}

\ifdefempty{\videocoverpath}{
{\small 请在生成笔记时填入视频封面图的本地路径。\par}
}{
\includegraphics[width=0.82\textwidth,height=0.45\textheight,keepaspectratio]{\videocoverpath}\par
}

\vfill
\begin{tcolorbox}[width=0.9\textwidth, colback=black!2!white, colframe=black!60, sharp corners]
\textbf{视频作者/频道}：\videochannel\par
\textbf{发布日期}：\videopublishdate\par
\textbf{视频时长}：\videoduration\par
\textbf{视频链接}：\href{\videourl}{\nolinkurl{\videourl}}
\end{tcolorbox}
\end{titlepage}

\tableofcontents
\newpage

\begin{figure}[H]
\centering
\includegraphics[width=0.78\textwidth]{overview_cover.jpg}
\caption{课程官方材料中的代表性页面，用于帮助读者在进入本讲正文前建立整体上下文。}
\end{figure}
\section{本讲定位与核心观点}

这是课程最后一讲，也是在整门课里最明确讨论“长程规划（long-horizon planning）”和“语言（language）”关系的一讲。讲者一开始就把观点说得很直接：人类行为本来就是层级式的，机器人如果想真正处理复杂任务，就不能只看动作本身，还必须把动作背后的目标、计划和语义纳入进来。

\subsection{ACTION = MOVEMENT + GOAL}

\begin{importantbox}{本讲最重要的公式}
\[
\text{ACTION} = \text{MOVEMENT} + \text{GOAL}
\]
动作（action）既是物理运动（movement），也是目标和意图（goal）的组合。只学其中一半，都会丢掉另一半的信息。
\end{importantbox}

讲者用 Harry Potter 的例子说明这个差异：书更擅长传达角色目标和意图，电影更擅长展示物理动作和可见运动。机器人学习需要同时理解这两层信息。

\subsection{为什么这与 LLM 相关}

大语言模型（large language models, LLMs）擅长抽象目标、常识和高层语义，但这并不意味着它们天然擅长低层运动控制。课程的态度很明确：语言模型很重要，但它不是“自动得到世界模型”的万能钥匙。

\subsection{本章小结}

本讲的起点是一个层级式世界观：语言、计划、目标、动作和物理运动并不是同一种东西。机器人系统必须同时处理这些层次。

\section{SHAKEY 与符号规划}

\subsection{早期规划系统}

SHAKEY（Nilsson, 1984）是机器人规划史上的经典早期系统之一。课程把它作为一个非常直观的起点：如果目标是堆叠方块，那么机器人可以先把“拿起”“放下”等动作形式化成逻辑算子（logical operators）。

\subsection{算子由什么构成}

每个算子（operator）都可以看成由三部分组成：
\begin{itemize}
  \item pre-condition list：执行动作前必须满足的条件；
  \item delete list：动作执行后不再成立的事实；
  \item add list：动作执行后新成立的事实。
\end{itemize}

这套表示法的好处在于，规划问题可以直接变成逻辑推理：从初始状态出发，找出一串算子，使得状态最终到达目标状态。

\subsection{本章小结}

SHAKEY 的历史意义在于，它把“规划”正式变成了一个可以用逻辑和算子讨论的问题。但它也非常依赖人工定义的动作和条件。

\section{PDDL：把规划写成正式语言}

\subsection{PDDL 是什么}

规划域定义语言（Planning Domain Definition Language, PDDL）是一种用来描述规划域、计划和目标的正式形式。讲义把它拆成四个关键部分：
\begin{enumerate}
  \item object types：对象类型；
  \item predicates：描述对象属性和关系的谓词；
  \item actions：动作本身；
  \item parameters / preconditions / effects：动作参数、前提和效果。
\end{enumerate}

\subsection{PDDL 的边界}

PDDL 很适合做符号规划（symbolic planning），但它是抽象层的语言，并不直接考虑动作如何在物理世界里执行。换句话说，PDDL 描述“应该做什么”，但不负责回答“如何在真实世界把它做出来”。

\subsection{本章小结}

PDDL 是经典规划理论的重要工具，但它的抽象程度很高。只靠 PDDL，还不能直接解决机器人在真实世界里的接触、碰撞和控制问题。

\section{TAMP：任务与运动规划}

\subsection{为什么需要 TAMP}

任务与运动规划（Task And Motion Planning, TAMP）试图把高层符号规划和低层运动规划结合起来。高层规划决定做什么，低层规划决定怎么做，二者联合起来才可能在真实机器人系统里工作。

\subsection{TAMP 与 PDDL 的区别}

PDDL 主要关心符号层的状态和动作，而 TAMP 进一步把机器人自身的物理构型、运动约束和碰撞避免也纳入规划过程。讲义里的核心意思是：高层规划可以给低层运动规划提供目标和约束，但低层是否能安全到达，仍然要由物理运动规划来保证。

\subsection{本章小结}

TAMP 是从“纯符号”走向“符号 + 物理执行”的关键一步。它说明真正的机器人规划不能只在逻辑层发生。

\section{语言模型与算子}

\subsection{LLM 可以帮什么忙}

讲义指出，现代 LLM 往往隐含了大量常识知识，因此可以用来帮助推断谓词和算子，减轻机器人研究者手工穷举 domain 的负担。换句话说，LLM 可能帮助我们更快地定义“做什么”和“什么时候能做”。

\subsection{但它仍然有局限}

课程同时提醒，LLM 仍然很难稳定地产生逻辑一致的计划。为了缓解这一点，一些工作开始结合程序合成（program synthesis）或其他结构化方法。这里的核心判断不是“语言模型没用”，而是“语言模型更像高层规划助手，而不是完整控制器”。

\subsection{本章小结}

LLM 最适合帮助机器人系统处理高层抽象：目标、约束、谓词和任务结构。它们并不自动等价于可靠的执行器。

\section{语言不擅长什么}

\subsection{空间认知（spatial cognition）}

课程后半部分转入一个很重要的反例：语言模型和 VLMs 在空间推理上并不强。
讲义引用了多个 benchmark 研究，指出现代模型在空间认知任务上普遍低于人类。

这类 benchmark 包括 mental rotation、shape matching、line orientation judgments 等。这些任务本质上要求模型精确地处理几何关系，而不是只给出高层语义解释。

\subsection{为什么低层动作不适合用语言承载}

讲义用“爬楼梯”的例子说明：如果真的想精确描述一串低层动作，语言会变得非常冗长，而且效率很低。相比之下，视频或模仿学习更适合传递动作级别的信息。

\subsection{本章小结}

语言并不是无所不能的表示形式。它适合高层抽象，不适合高精度空间关系，也不适合细粒度动作序列的直接传输。

\section{总结与延伸}

这一讲最后收束得很清楚：语言在机器人里很有价值，但它的价值在于表达抽象目标、环境属性和计划结构，而不是代替运动控制。

如果把整门课串起来看，可以得到一个非常稳定的系统观：
\begin{itemize}
  \item 监督学习擅长从示范里学局部动作；
  \item 局部策略擅长处理接触丰富的操作；
  \item 语言和符号规划擅长组织高层结构；
  \item 运动规划和控制负责把计划落到物理世界。
\end{itemize}

这也意味着，真正有希望的机器人系统，往往不是“一个模型统治全部”，而是层级化、模块化、互相协作的系统。

\subsection{本章小结}

本讲的终点不是“语言能解决机器人问题”，而是“语言能解决机器人问题中的高层部分；低层部分仍然需要几何、控制和动作学习”。

\end{document}
